\makeatletter
\def\input@path{{./}{../}{../../}{../../../}{../../../../}{preamble/}{../preamble/}{../../preamble/}{../../../preamble/}}
\makeatother

\input{preamble_exams.tex}
\usepackage{needspace}
\usepackage{titling}
\setlength{\droptitle}{-2em}
\pretitle{\par\begingroup\centering\LARGE}
\posttitle{\par\endgroup}
\preauthor{}\postauthor{}
\predate{}\postdate{}
\maketitlehookd{\vspace{-1em}}
\linespread{1.12}
\setlength{\parskip}{5pt}
\newgeometry{margin=0.65in,top=0.55in,bottom=0.65in}

\title{Grade 11 Term 2 Home Practice -- C5\\\large A population proportion\\[0.25em]
	\normalsize Global Economics -- 2026--2027}
\author{}
\date{}

\begin{document}
	\maketitle

	\setcounter{tocdepth}{1}
	\setcounter{secnumdepth}{2}
	\phantomsection
	\label{toc}
	\tableofcontents

	\section*{Evaluated criteria}
	\begin{enumerate}[label=\textbf{C\arabic*},start=5]
		\item Constructs confidence intervals for a population proportion using sample proportions, applying them to contexts such as poverty rates, defaulting borrowers, or successful investments.
	\end{enumerate}

	\textbf{How to practise.} Try each task before opening its blue solution link. The arrows return to the contents, problem list or current problem.

	\section{Worked Solution Examples}
	\label{sec:worked-solution-examples}

	\SectionProblemsTableOfContents{
		\SectionProblemLink{prob:c5-1}{First-year microloan defaults (C5)} \\
		\SectionProblemLink{prob:c5-2}{Export shipment delays (C5)}
	}

	% ----------------------------------------------------------------
	\subsection{First-year microloan defaults (C5)}
	\label{prob:c5-1}

	A lender takes a simple random sample without replacement of 400 loans from 8,000 small-business loans issued in the same year. Every sampled loan has a complete first-year record. Of these loans, 50 defaulted within the first year. For a two-sided 98\% normal interval use $z^*=2.326$.

	\QuestionsTasks[sec:worked-solution-examples]{
		{5}/{sol:c5-1-a}/
		{Define $p$, compute $\hat p$, and check approximate independence and the two outcome counts for a one-proportion $z$ interval.},
		{5}/{sol:c5-1-b}/
		{Compute the estimated standard error, margin of error and 98\% interval. Interpret the interval and compare it with a 20\% default benchmark.}
	}

	\clearpage

	\Needspace{5\baselineskip}
	\phantomsection
	\label{sol:c5-1-a}
	\subsubsection*{C5 -- Identify the parameter and check the procedure}

	Let $p$ be the fraction of all 8,000 loans in the cohort that defaulted within the first year. Using the Term 1 relative frequency idea,
	\[
	\hat p=\frac{x}{n}=\frac{50}{400}=0.125.
	\]
	The simple random sample covers $400/8{,}000=5\%$ of the cohort, within the 10\% guideline for approximate independence. The observed counts are $n\hat p=50$ defaults and $n(1-\hat p)=350$ non-defaults. Both are at least 10, supporting the normal approximation for this interval.

	\secInicio[prob:c5-1][sec:worked-solution-examples]

	\Needspace{5\baselineskip}
	\phantomsection
	\label{sol:c5-1-b}
	\subsubsection*{C5 -- Build and interpret the default-rate interval}

	\[
	\widehat{SE}=\sqrt{\frac{\hat p(1-\hat p)}{n}}
	=\sqrt{\frac{0.125(0.875)}{400}}\approx0.01654,
	\qquad E=2.326\,\widehat{SE}\approx0.03846.
	\]
	Using the unrounded intermediate values, $\hat p\pm E$ gives approximately $[0.08654,0.16346]$. We are 98\% confident that about 8.65\% to 16.35\% of this loan cohort defaulted within its first year. The entire interval lies below 20\%, providing evidence that the first-year default rate for this cohort was below the benchmark.

	\secInicio[prob:c5-1][sec:worked-solution-examples]

	\clearpage

	% ----------------------------------------------------------------
	\subsection{Export shipment delays (C5)}
	\label{prob:c5-2}

	An export council draws a simple random sample of 150 different shipments from 3,000 shipments dispatched this season. Of these, 39 arrived at least one day late. Use $z^*=1.96$ for a two-sided 95\% interval.

	\QuestionsTasks[sec:worked-solution-examples]{
		{5}/{sol:c5-2-a}/
		{Define the population proportion, check the conditions, and calculate a 95\% one-proportion $z$ interval for the delay rate.},
		{5}/{sol:c5-2-b}/
		{Express the point estimate and margin of error in percentage form with appropriate units. Interpret the interval and assess whether it establishes that the delay rate is above 25\%.}
	}

	\Needspace{5\baselineskip}
	\phantomsection
	\label{sol:c5-2-a}
	\subsubsection*{C5 -- Construct a 95 percent interval}

	Here $p$ is the fraction of all 3,000 shipments delayed by at least one day. The random sample covers $150/3{,}000=5\%$ of the population. There are 39 delayed shipments and $150-39=111$ other shipments, so both counts meet the threshold of 10. With $\hat p=39/150=0.26$,
	\[
	\widehat{SE}=\sqrt{\frac{0.26(0.74)}{150}}\approx0.03581,\qquad
	E=1.96\,\widehat{SE}\approx0.07020.
	\]
	Using the unrounded margin, the interval is approximately $[0.18980,0.33020]$.

	\secInicio[prob:c5-2][sec:worked-solution-examples]

	\Needspace{5\baselineskip}
	\phantomsection
	\label{sol:c5-2-b}
	\subsubsection*{C5 -- Report the estimate and its precision}

	The point estimate is 26.0\%, with a margin of about 7.02 percentage points. We are 95\% confident that about 19.0\% to 33.0\% of all shipments this season were delayed by at least one day. Since 25\% and values below it are within the interval, the data leave open whether the population delay rate exceeds 25\%.

	\secInicio[prob:c5-2][sec:worked-solution-examples]


\end{document}
